\section{Row Sums and Transpose: a Chain of Specific Lemmas}\label{sec:rows}

\subsection{The program}
\begin{lstlisting}
method main(A: array<int, array<int, int>>, m: int, n: int) returns (s: array<int, int>)
  requires 0 <= m && 0 <= n
  ensures forall i: int :: 0 <= i < m ==>
    s[i] == sum(selects(A[i], butlast(counter(0, 1, natural(n)))))
{
  var i: int;
  var j: int;
  var t: int;
  i := 0;
  while i < m
  {
    t := 0;
    j := 0;
    while j < n
    {
      t := t + A[i][j];
      j := j + 1;
    }
    s := s[i := t];
    i := i + 1;
  }
}
\end{lstlisting}
The inner loop sums one row and the outer loop stores the sum at the position of its iteration; the
transpose (\texttt{transpose.dfy}) is the same shape with an array-valued inner result, the stores of
column $k$ into $B$. Both are run with \texttt{-include teq:nested}, the module of the fragments of inner loops.

\subsection{The trace equation system}
With the module, the goal of the program alone is
\begin{align}
  K &= I = \cnt(0, 1, n') \qquad \intf(n') = m \label{eq:rs0}\\
  &\al{\bl{K}, \bl{I}, \seq{N}, \seq{K}'}{k, i, nn, K' \mapsto K' = \cnt(0, 1, nn)} \nonumber\\
  &\al{\bl{K}, \seq{N}, \seq{J}'}{k, nn, J \mapsto J = \cnt(0, 1, nn)} \label{eq:rs1}\\
  &\seq{N} = \bl{\mathsf{consts}(n', \natf(n))} \qquad \mathsf{consts\_prefix}(n', n, \seq{J}) \label{eq:rs2}\\
  &\al{\bl{K}, \bl{I}, \seq{T}', \seq{J}'}{k, i, T, J \mapsto T = \sums(0, \mathsf{selects}(A[k], \bl{J}))} \label{eq:rs3}\\
  &\al{\bl{K}, \bl{I}, \seq{T}, \seq{T}'}{k, i, t, T \mapsto t = \cur(T)} \label{eq:rs4}\\
  &\al{\bl{K}, \bl{I}, \seq{J}}{k, i, j \mapsto n \le j} \qquad \ldots \nonumber\\
  &\Longrightarrow\ 0 \le i < m \Rightarrow \mathsf{earlier}(\natf(i), \seq{T}) = \mathsf{sum}(\mathsf{selects}(A[i], \bl{\cnt(0, 1, \natf(n))})) \nonumber
\end{align}
The module has resolved the inner counter \eqref{eq:rs1}, the exit (abstracted to $\mathsf{consts\_prefix}$)
and the numbers of inner iterations \eqref{eq:rs2}: $\seq{N}$ is the constant sequence $n$. The sum of
row $k$ is the trace \eqref{eq:rs3} over the inner counter, its result the fragment \eqref{eq:rs4}, and the
postcondition reads the stored results at position $i$. What is left is to carry $\seq{N}$ into the counter,
the counter into the sum, the sum into the result, and the results into the read: four substitutions
between fragments over different dependency lists, which the simplifier does not make by itself.

\subsection{The chain}
The first version of this file stated the substitutions as general lemmas: a fragment about the inner
counter $J$ with a higher-order body $t = F(k, J)$ is stated about $\cnt(0, 1, nn)$, and a read
$\mathsf{earlier}(\natf(i), ts)$ is the body of the $\mathsf{all}$ over $ts$ at $i$. Both have a general
left-hand side, matching every fragment, respectively every read, together with an $\mathsf{all}$ side
condition over the lists of the left-hand side; matching the condition builds the simplifier's context of
the fact's combinator, where the reads of the dependencies are rewritten, on which the rule matches again,
one level deeper, until the step budget is exhausted (finding 19 of \texttt{NOTES.md}, the cycle that
\texttt{TODO.md} describes for \lemma{sum\_invs\_const}). The runs did not terminate, which the runners of
the time reported as proofs, since a killed run prints no goal.

The version that verifies has one specific lemma per step, each with the fragment it depends on as a side
condition over a fresh list variable, identified by an equation (Section~\ref{sec:cock}); $C_n$ abbreviates
the counter $\bl{\cnt(0, 1, \natf(n))}$:
\begin{align*}
  &\lemma{invs\_counter\_consts}:\quad 0 \le n\\
  &\quad\Longrightarrow\ \al{ks, \bl{\mathsf{consts}(m, \natf(n))}, Js}{k, nn, J \mapsto J = \cnt(0, 1, nn)}\\
  &\qquad\Longleftrightarrow\ \al{ks, Js}{k, J \mapsto \bl{J} = C_n}\\[2pt]
  &\lemma{invs\_sums\_selects\_counter}:\quad \al{ks, Js_2}{k, J \mapsto \bl{J} = C_n}\ \wedge\ Js_2 = Js\\
  &\quad\Longrightarrow\ \al{ks, Ts, Js}{k, T, J \mapsto T = \sums(0, \mathsf{selects}(F(k), \bl{J}))}\\
  &\qquad\Longleftrightarrow\ \al{ks, Ts}{k, T \mapsto T = \sums(0, \mathsf{selects}(F(k), C_n))}\\[2pt]
  &\lemma{invs\_current\_def}:\quad \al{ks, Ts_2}{k, T \mapsto T = G(k)}\ \wedge\ Ts_2 = Ts\\
  &\quad\Longrightarrow\ \al{ks, ts, Ts}{k, t, T \mapsto t = \cur(T)}\\
  &\qquad\Longleftrightarrow\ \al{ks, ts}{k, t \mapsto t = \cur(G(k))}\\[2pt]
  &\lemma{invs\_maps\_sum}:\quad \al{ks, ts}{k, t \mapsto t = \mathsf{sum}(F(k))}\\
  &\qquad\Longleftrightarrow\ ts = \mapc{ks}{k \mapsto \mathsf{sum}(F(k))}\\[2pt]
  &\lemma{earlier\_maps\_counter}:\quad 0 \le i < \intf(m)\\
  &\quad\Longrightarrow\ \mathsf{earlier}(\natf(i), \mapc{\bl{\cnt(0, 1, m)}}{k \mapsto F(k)}) = F(i)
\end{align*}
The first matches \eqref{eq:rs1} with the dependency $\bl{\mathsf{consts}(m, \natf(n))}$ as a pattern,
which is matched against $\seq{N}$ through its equation \eqref{eq:rs2} in the context (a variant with
$\seq{N}$ as a variable and the equation as a side condition is in the file as well). The second rewrites
\eqref{eq:rs3} to a sum over the counter $\bl{\cnt(0, 1, \natf(n))}$, with $F := \lambda k.\ A[k]$. The third
rewrites \eqref{eq:rs4} with $G$ the sum of row $k$, after which the database makes $\cur(\sums(0, xs))$ into
$\mathsf{sum}(xs)$; the fourth turns the pointwise fragment of the results
into a map (specific to sums, so that the fragments $j = \cur(J)$ of counters are left to the module); the
fifth reads the map at $i$, with $m := n'$ and $0 \le i < \intf(n')$ decided from $0 \le i < m$ and
\eqref{eq:rs0}. Every left-hand side is a fragment with a function symbol on the side of its bound variable
or a read of a $\mathsf{map}$, and every $\mathsf{all}$ condition is over a fresh list: nothing matches what
the rewriting of a read in a combinator context produces, and the runs take under twenty seconds.

For the transpose the chain is the same with the inner stores: the trace of stores over $\bl{J}$ is folded
into $\mathsf{column}(h, A, k, \natf(n))$ once $J = \cnt(0, 1, \natf(n))$ is known, the final values of the
columns are the fragment $b = \cur(\mathsf{column}(h, A, k, \natf(n)))$ over the previous value $h$, then
\lemma{current\_results\_steps} (Section~\ref{sec:nsum}) makes the outer trace the trace of the columns, and
\lemma{current\_transpose} reads the transposed element.

\subsection{The idea}
An inner result stored at the position of the outer iteration is the pattern of the paper's nested loops,
and with the module it needs no new concept, only discipline: every step is a lemma whose left-hand side is
specific and whose side condition is a fragment over a fresh list, matched against the facts. The general
substitution lemmas that one would write first are exactly the rules that the cycle mechanism of
\texttt{TODO.md} forbids; the chain makes the same substitutions one fragment at a time.
